\documentclass[11pt]{article}

% Change "review" to "final" to generate the final (sometimes called camera-ready) version.
% Change to "preprint" to generate a non-anonymous version with page numbers.
\usepackage[preprint]{acl}

% Standard package includes
\usepackage{times}
\usepackage{latexsym}
\usepackage[]{hyperref}

% For proper rendering and hyphenation of words containing Latin characters (including in bib files)
\usepackage[T1]{fontenc}
% For Vietnamese characters
% \usepackage[T5]{fontenc}
% See https://www.latex-project.org/help/documentation/encguide.pdf for other character sets

% This assumes your files are encoded as UTF8
\usepackage[utf8]{inputenc}

% This is not strictly necessary, and may be commented out,
% but it will improve the layout of the manuscript,
% and will typically save some space.
\usepackage{microtype}

% This is also not strictly necessary, and may be commented out.
% However, it will improve the aesthetics of text in
% the typewriter font.
\usepackage{inconsolata}

%Including images in your LaTeX document requires adding
%additional package(s)
\usepackage{graphicx}

% If the title and author information does not fit in the area allocated, uncomment the following
%
%\setlength\titlebox{<dim>}
%
% and set <dim> to something 5cm or larger.

\title{How Educational Differences are Reflected in Reddit Posts}

% Author information can be set in various styles:
% For several authors from the same institution:
% \author{Author 1 \and ... \and Author n \\
%         Address line \\ ... \\ Address line}
% if the names do not fit well on one line use
%         Author 1 \\ {\bf Author 2} \\ ... \\ {\bf Author n} \\
% For authors from different institutions:
% \author{Author 1 \\ Address line \\  ... \\ Address line
%         \And  ... \And
%         Author n \\ Address line \\ ... \\ Address line}
% To start a separate ``row'' of authors use \AND, as in
% \author{Author 1 \\ Address line \\  ... \\ Address line
%         \AND
%         Author 2 \\ Address line \\ ... \\ Address line \And
%         Author 3 \\ Address line \\ ... \\ Address line}

\author{Maxwell Champlin \\
  3859119 / Maxc1102@gmail.com \\}


\begin{document}
\maketitle
\begin{abstract}
In this project we sought to understand the linguistic feature differences by education of reddit posters. As such we used 5000 posts from the r/PhD and r/Teenagers subreddit to be our main dataset for feature differences. As well, we wished to be able to classify posts by subreddit, for this we used 1000 posts. We sought to look at as many linguistic features as possible, as well as specialized educational tools like the Brysbert age of word acquisition and the Flesch-Kincaid reading ease and grade level scores. We found that the key differences between the two subreddits did not lay in general linguistic features, and that features like the amount of punctuation that we had expected to show higher levels of education found that teenagers used more than posters in PhD. Age of acquisition and reading grade level scores showed differences in the direction we had expected, however, the reading level and age of word acquisition of posts for both subreddits were well below our expectations. For our classifier we found that classifying the subreddits was relatively easy when looking at the total samples, as the words used between them were often very different. These results highlight the fact that expectations of educational differences showing up in the linguistic features of text is not likely, at least on informal sources, and that education is best found by looking at lexical selection. 
\end{abstract}

\section{Introduction}


Language is an innate part of the human cognition system. Perhaps it is through Chomsky's proposed genetic faculty for language, or perhaps not. What we can be certain of is that reading and writing are not natural skills that humans are imbued with, instead they are skills that take many years to learn and are require artificial constructions like spelling and punctuation to convey an intended meaning, this is not a crutch that spoken language needs. Using subreddit data from r/Teenagers and r/PhD, we hope to be able to examine the linguistic differences between the two groups and be able to examine how length of time in education plays a roll in various linguistic features. In particular we will be examining the overall reading grade level and reading ease of the posts using the Flesch-Kincaid reading ease and grade level calculations. This will allow us to get an overview of how these two communities differ lexically. We will also be using a number of features, like the amount of punctuation, number of unique nouns per post, and more to see how the communities vary and if our predictions of more formal tools like punctuation being something that is mostly used by those with more education. 

%RQ1: how does educational differnces show up in overall reading grade level of writing
%RQ2: do more educated people use more variated vocabularly in reddit posts
%RQ3: Do more educated people use more punctuation and other more 'formal' writing tools 

\section{Data}
In this project the dataset used came from two subreddits: r/Teenagers and r/PhD. The data used was a subsample of 5000 posts, which was then filtered to only include those over 10 and less than 500 words. In total we had about 2696 posts for the Teenager subreddit and about 3286 for the PhD subreddit. The choice for 5000 posts which was then filtered was made considering the statistical significance of our set over the 1 million or so posts in the teenagers dataset. As well, the size was chosen to reflect the limitations on computing power as each additional post added about a second of processing time for the machine. 


For the classifier we created, the entirety of the PhD subreddit was loaded, while only a subset of the Teenager subreddit was added to our classifier. There were 10,001 utterances for each dataset, there were 2700 speakers in the PhD dataset, and 1699 for the teenager dataset. This as well as the smaller number of conversations in the Teenager data, informs us that there are more long threads of people replying to each other in the Teenagers subreddit. 



%data for the phd subreddit
In the PhD subreddit our median number of words (post filtering) per post is 78, this is higher than the teenagers who are at 65, with standard  deviations of 95 and 94 respectively. This highlights one of the most interesting parts of our data, that for the majority of our features the overall distributions are surprisingly flat. There are key differences between the two subreddits, which we will see later, but for the most part the linguistic features are very normal and do not speak to huge differences. However, in our reading ease and grade level, as well as age of word acquisition statistics there are clear differences, which we will see in further sections. 

\section{Features}
Describe the set of features you are going to use. Give examples for features that are a bit less intuitive to interpret. Repeat a hypothesis that you have regarding the relationship between certain features and certain empathy types or how they differ between communities or speakers. 

For the research questions we will use a mixture of pre-built and manual features from LFTK to analyze the data. In the area of how educational differences show up in the posts of r/PhD posters or r/Teenager posters we looked at the prebuilt LFTK Flesch-Kincaid Grade Levels\footnote{The Grade level used are American Grade Levels. The score is calculated by dividing the total words by total sentences and adding it to the total syllables over total words. The exact equation can be found \hyperlink{https://web.archive.org/web/20201210212716/https://apps.dtic.mil/sti/pdfs/ADA006655.pdf}{Here} }, the Flesch-Kincaid Reading Ease\footnote{The Reading Ease score is calculated by dividing the total words by total sentences, and then minus the total syllables over total words. There are a number of scalars in the calculation to be able to get the 0-100 scale. Further information \hyperlink{https://web.archive.org/web/20160712094308/http://www.mang.canterbury.ac.nz/writing_guide/writing/flesch.shtml}{Here} }, and the Brysbert age of acquisition of words, and of sentences\footnote{The age of acquisition is determined through surveys and currently contains about 44,000 words.}.



The specialized metrics for grade level and reading ease were chosen due to them offering a clear measurement of lexical differences between the two subreddits. Particularly the age of acquisition metrics offer a clear look at the lexical differences with a more robust methodology than the reading ease and grade level which, while useful, only concern themselves with the number of syllables per word, and could be less accurate than expected due to proper nouns or cases where words are not properly spaced. The use of the reading ease and grade level statistics is that they allow us to see the overall differences in word length, which while not a perfect representation of lexical choices allows us to see a rough idea of what the overall lexical differences are by word length. 

As well, as LFTK is a powerful tool, we are also looking at a full suite of syntactic features in the data. The key features that we intend to look at based on our hypotheses of what linguistic features correlate with education levels, we expect that PhD will have more punctuation, as that is often associated with more formal writing. As well we expect the number of unique nouns to be higher for PhD as we expect that more educated people will use more unique nouns than teenagers will. Since we have the ability to look at a large number of features, we will report later surprise results from our feature variation examination. 

\section{Methodology}
%Describe the steps you are going to perform in order to retrieve results. 
%(1) pre-processing (do you convert the dataset to a Corpus Object? Do you do any cleaning of the data?)
%(2) feature extraction: which tool(s) do you use?
%3) Experimental setup for prediction task: What are your data splits? Best practice for all experiments (now and in the future): always use several (different) folds (at least 3) for training + evaluation. Report the average and standard deviation of performance results. If your target label is imbalanced, report F1 Macro Score and not accuracy. Also report the result for the positive and negative class (so you know whether you are overpredicting or undergeneralizing). 
%(4) For the feature analysis (either correlation with empathy or variation between communities), what methods are you using? E.g. plots, statistical model, or maybe you are calculating correlations? Identifying discriminating words based on frequency? Feature importance in the classification?



To begin with our preprocessing we will break it down into two sections, as the classification and linguistic feature analysis used different preprocessing processes. We will begin by looking at the linguistic analysis. The first step was using the Cornell subreddit collections, downloading the entirety of the two communities' subreddits. From there they are loaded as a corpus file, and tokenized to create the field $'num_tokens'$ which is needed for further steps. From there the full corpora are put through the spacy English core small model via an LFTK extractor, the data (our LFTK features and values) is then put into a CSV file. From there the CSV file is then read into a dataframe, which is concatenated with the original Corpus objects. After the concatenation the posts with either above 500 or below 10 tokens are removed from the data\footnote{This is the main reason that only 5000 utterances were considered, this is a very inefficient way of working with the data, and unfortunately I ran out of time to rework code that already worked on statistically significant amounts of data. }. From there, we use the data to create charts and analyses using a number of statistical tools. We will see more about this in a further section.



For the classifier, data was input directly from the corpora, without filtering, this was due to the effects of short and long messages on the classifier being minimal compared to the same impacts on the linguistic feature analysis, as some features are per post and long or short posts create outliers the unfortunate ability to influence overall data. We had 10,000 utterances for both subreddits, but over 2900 conversations for the PhD subreddit, and only 575 for the Teenager subreddit. It is possible that the increase of crosstalk (i.e. posters talking to each other in comments instead of responding to main post) may have had a small influence, but we felt it was a reflection of the communities and was best left in the data. After that the utterance data was put into dataframes for both communities. After this step, we merged the datasets to create a single dataframe with labels for the correct subreddit. We then break the data into test and training sets. The training and testing sets contain both an input and the output–since we will be predicting the subreddit from the utterance, the utterance is the input and the label is the output. The utterances are in text form and must be converted to a numerical form. Two methods were used for this–bag of words and TF-IDF, including an attempt to improve accuracy through n-gram features.



\section{Results}
In the following we will examine the results of the linguistic feature analysis and the classifier.  Let us begin with the Average age of acquisition of words per sentence, as we can see in Figure \ref{Figure 1}. 

\begin{figure}[h!]
\includegraphics[height=4cm, width = 6cm]{a_kup_ps.png}
\caption{Average Kuperman age of acquisition per sentence}
\label{Figure 1}
\end{figure}
We can clearly see that the age of acquisition is higher for the PhD community with a mean of 81.6, and a standard deviation of 35.5, while the Teenager community having a mean of 64.8 with a standard deviation of 29.4. The high standard deviation will be seen throughout the data, this is caused by a wide variation in the type of post, with many posts being personal in nature, while some are highly technical, this can be seen in the larger distribution towards the high end for the PhD community. The higher prevalence in personal stories in r/Teenagers may also explain the distribution of pronouns, which we can see in Figure \ref{figure 2}

\begin{figure}[h!]
\includegraphics[height=3cm, width = 7cm]{n_pron.png}

\caption{number of pronouns}
\label{figure 2}
\end{figure}

With a standard deviation of around 0.05 for both communities, we can see a very large difference in the number of pronouns. Why this is, is not yet established, but it is not an expected outcome for what we imagined when comparing education. As such it points to another confounding variable for analyses of linguistic features using Reddit. The composition of the type of post (personal stories, tech help, general forum posting, etc) is not at all stable across subreddits. Future work could seek to understand what the distribution of conversational types is across subreddits. 

\begin{figure}[h!]
\includegraphics[height=3cm, width = 7cm]{num_tokens.png}

\caption{number of tokens}
\label{figure 3}
\end{figure}

As well, as we can see in \ref{figure 3} above, the number of tokens (words). Here the difference is relatively small, with standard deviation of 95 for both subreddits, it is clear that the posts are overall about the same, but we can see that PhD is much less normal than Teenagers. The thicker lune shows us more posts from PhD nearing the 500 word filter for our data. As well, Teenagers has a notably more normal. 
\begin{figure}[h!]
\includegraphics[height=3cm, width = 7cm]{n_punct.png}

\caption{number of punctuation marks}
\label{figure 4}
\end{figure}
Contrary to our initial hypothesis, the amount of punctuation between the subreddits is weighted towards Teenagers instead of PhD, as seen in Figure \ref{figure 4}. This is not something that we have clear explanations for, but with a standard deviation of 0.09 for Teenagers and 0.05 for PhD, this is a clear trend in the data. 

\begin{figure}[h!]
\includegraphics[height=3cm, width = 7cm]{fkgl.png}

\caption{Flesch-Kincaid reading grade level}
 \label{figure 5}
 \end{figure}
 
 In Figure \ref{figure 5} above we can see that the reading grade level for PhD is higher than that of Teenagers. This is of course to be expected, as we hypothesize that the more educated speakers will have a different lexicon than less educated speakers. One might note that the overall level as calculated by this metric is rather low for both PhD and Teenagers, who both appear to be writing well below what their respective levels of education would suggest. This again leads us to the issue of the type of conversation, personal stories and other such posts do not highlight lexically the educational differences. 
 
 
 As well, in Figure \ref{figure 6}, we can see very similar data with the reading ease metric. Which allows us to see the data without the scaling of the American school system 
 \begin{figure}[h!]
\includegraphics[height=3cm, width = 7cm]{fkre.png}
\caption{Flesch-Kincaid reading Ease}
\label{figure 6}
\end{figure}


\subsection{Predicting source community}
For our classification task we ran three different models, a Bag of Words model, a TF-IDF model, and a N-gram model. 
For the BoW model the entirety of the ‘input’ for the training data is made into a list of all of the unique words inside of it. And then, for each utterance, a vector is created to indicate which of those unique words appear in the utterance, along with how many times each of those words appears. This allows for a rough way for the model to distinguish between the communities, but there is no control for frequent words like \textit{the}. 
The n-gram is the same as the BoW but instead of the list of all unique words, it also includes unique pairs of words (since the n-gram range was set to (1, 2)). 
And our final model the, TF-IDF builds upon the BoW, as it keeps track of how many times each unique word appears in a given utterance, then factors in how often that word appears across the entire training dataset. By factoring in rarity, it punishes words like ‘the’ and ‘is’ for being too common. So for example, TF-IDF would treat the word ‘the’ the same as the word ‘thesis’ if they were 
equally common in a given utterance.

For the training, the classifiers observed the training dataset, then draws a hyperplane that maximizes distance between the two clusters on either side, to make a prediction it simply sees which side of the line the point falls on. 
The accuracy of the three models were all well above random chance, with the n-gram performing about the same as the Bag of Words model, while the TF-IDF model performed the best with about an 85\% accuracy over our test data. As well, with the models we saw the strongest correlations with lexical differences. Words like 'high school' and 'university' were key differentiators 


\begin{center}
\begin{tabular}{|| c |c | c | c || }
\hline
 Model & Bag of words & TF-IDF & n-gram   \\ 
 \hline
 \hline
 Accuracy & 0.818 & 0.84 & 0.81 \\  
 \hline

\end{tabular}
\end{center}



\section{Discussion}
%Reflect on your findings: did you see expected results? Did you see something unexpected? What data points were missed or miss-classified by the model(s) and why? What other features would be promising? What additional information would be useful to know or annotate to perform better? Did you find particular properties of the data that made classification more or less reliable (e.g. reddit posts are over a long time span, the way empathy is annotated is oversimplified ...)


Overall the linguistic features were not exactly what we had expected, while the classification went well, it showed us the general trend, we are looking at lexical differences instead of having clear linguistic feature differences, with a few exceptions. As well, we did not consider the tone and use of subreddits, and that will have biased our work away from comparing like to like posts, like questions vs questions or stories vs stories. This then, leaves us with results that are unexpected from our initial expectations, at least in part, and leave open a goodly amount of ground for further work in this area. Our classification work as well highlighted the key differences between the subreddits, as they are easily distinguished even by simple models, including an n-gram model, which tells us, even without looking directly at the words used, along with our reading ease and grade level metrics that we are looking at nearly a purely lexical difference between the two subreddits. 

\section{Conclusion}
Taking 5000, and 10000 posts from the Teenagers and PhD subreddits we sought to examine how linguistic differences through education present themselves on reddit. By using metrics like Flesch-Kincaid reading ease and grade level, as well as Brysbert age of acquisition of words, we had a number of specialized tools to examine the lexical differences between the two subreddits. As well, we had a whole suite of general linguistic features. 

We showed that there are not major linguistic feature differences between the Teenagers and the PhD subreddits, this is contrary to our initial expectations. Instead while there are differences they are mainly limited to the areas of post type (see higher use of pronouns in r/Teenagers), and the lexical variation is the key factor between them. 
\appendix


\section{Use of generative AI}
No generative AI was used in the writing of this paper. No AI was used in the writing of the Corpus Load and Filter or Analysis files. I cannot speak to the use of GenAI tools in the NLP\_A1 file.  

\subsection{Link to Github}

The code for this assignment can be accessed at this link: \url{https://github.com/maxchamplin/NLP_and_Identity_project1}


\end{document}
